\chapter{Convergence}
\label{chap:pet-convergence}

\newcommand{\Ric}{\operatorname{Ric}}
\newcommand{\Hess}{\operatorname{Hess}}
\newcommand{\vol}{\operatorname{vol}}
\newcommand{\diam}{\operatorname{diam}}
\newcommand{\inj}{\operatorname{inj}}
\newcommand{\sec}{\operatorname{sec}}
\newcommand{\Cap}{\operatorname{Cap}}
\newcommand{\Cov}{\operatorname{Cov}}
\newcommand{\dGH}{d_{G\text{-}H}}

\section{Gromov-Hausdorff Convergence}

\subsection{Hausdorff Versus Gromov Convergence}

\begin{definition}[Hausdorff distance]
  \label{def:pet-ch11-hausdorff-distance}
  \lean{PetersenLib.hausdorffDistance}
  \source{petersen:page-0409}
  Given a metric space $(X,|\cdot|)$ and subsets $A,B\subset X$, define
  $d(A,B)=\inf\{|ab|\mid a\in A,\,b\in B\}$,
  $B(A,\varepsilon)=\{x\in X\mid|xA|<\varepsilon\}$ (writing $|xA|=d(\{x\},A)$),
  and define the \emph{Hausdorff distance}
  \[
    d_H(A,B)=\inf\{\varepsilon\mid A\subset B(B,\varepsilon),\,B\subset B(A,\varepsilon)\}.
  \]
  Thus $d(A,B)$ is small when some points of $A,B$ are close, while
  $d_H(A,B)$ is small exactly when every point of $A$ is close to a point of
  $B$ and vice versa. The Hausdorff distance defines a metric on the compact
  subsets of $X$, and this collection is compact when $X$ is compact.
\end{definition}

\begin{convention}[compact and proper metric spaces]
  \label{conv:pet-ch11-compact-proper-spaces}
  \lean{PetersenLib.ProperMetricSpace}
  \source{petersen:page-0409}
  We restrict to compact or \emph{proper} metric spaces:
  spaces whose distance function is proper, i.e.\ all closed balls are
  compact. Proper metric spaces are automatically separable, complete, and
  locally compact.
\end{convention}

\begin{definition}[admissible metric, Gromov--Hausdorff distance]
  \label{def:pet-ch11-gromov-hausdorff-distance}
  \lean{PetersenLib.admissibleMetric,PetersenLib.gromovHausdorffDistance}
  \source{petersen:page-0409}
  \uses{def:pet-ch11-hausdorff-distance}
  For two metric spaces $X,Y$, call a metric on the disjoint union
  \emph{admissible} when it extends the given metrics on $X$ and on $Y$. The
  \emph{Gromov--Hausdorff distance} is
  \[
    \dGH(X,Y)=\inf\{d_H(X,Y)\mid\text{admissible metrics on }X\cup Y\}:
  \]
  one places a metric on $X\cup Y$ extending those on $X,Y$ so as to make $X$
  and $Y$ as Hausdorff-close as possible — equivalently, one tries to define
  distances between points of $X$ and $Y$ without violating the triangle
  inequality.
\end{definition}

\begin{example}[distance to a point]
  \label{ex:pet-ch11-gh-distance-radius}
  \lean{PetersenLib.ghDistance_le_radius}
  \source{petersen:page-0409}
  \uses{def:pet-ch11-gromov-hausdorff-distance}
  \dcref{ch11:11.1.1}
  When $Y$ consists of one point,
  $\dGH(X,Y)\le\operatorname{rad}X=\inf_{y}\sup_{x\in X}|xy|$, the radius of
  the smallest ball covering $X$.
\end{example}

\begin{example}[diameter bound]
  \label{ex:pet-ch11-gh-distance-diameter-bound}
  \lean{PetersenLib.ghDistance_le_halfDiam}
  \source{petersen:page-0409}
  \uses{def:pet-ch11-gromov-hausdorff-distance}
  \dcref{ch11:11.1.2}
  If both diameters are at most $D$, assign cross-distance $D/2$ to every
  pair in $X\times Y$; this admissible metric gives $\dGH(X,Y)\le D/2$.
\end{example}

\begin{proposition}[distance zero implies isometric]
  \label{prop:pet-ch11-gh-distance-zero-isometric}
  \lean{PetersenLib.ghDistance_eq_zero_iff_isometric}
  \source{petersen:page-0410}
  \uses{def:pet-ch11-gromov-hausdorff-distance}
  \dcref{ch11:11.1.3}
  Compact metric spaces at zero Gromov--Hausdorff distance are isometric:
  $\dGH(X,Y)=0$ implies that $X$ and $Y$ are isometric.
\end{proposition}
\begin{proof}
  \uses{prop:pet-ch11-gh-distance-zero-isometric}
  \sketch
  Choose admissible metrics with \(d_H(X,Y)<i^{-1}\), and maps \(I_i:X\to Y\), \(J_i:Y\to X\) moving each point by at most \(i^{-1}\). Then
  \[
    |I_i(x_1)I_i(x_2)|\le |x_1x_2|+2i^{-1},\quad
    |J_i(y_1)J_i(y_2)|\le |y_1y_2|+2i^{-1},
  \]
  while \(J_iI_i\) and \(I_iJ_i\) move points by at most \(2i^{-1}\). Compactness and diagonal subsequences on countable dense sets give distance-nonincreasing limits \(I\) and \(J\). The last estimates make them inverse, hence isometries.
\end{proof}


\begin{remark}[$(\mathcal{M},\dGH)$ as a metric space]
  \label{rem:pet-ch11-gh-pseudo-metric-quotient}
  \lean{PetersenLib.compactMetricSpaceQuotient}
  \source{petersen:page-0410}
  \uses{prop:pet-ch11-gh-distance-zero-isometric}
  Symmetry and the triangle inequality make $\dGH$ a pseudometric on the
  collection $\mathcal{M}$ of compact metric
  spaces; by \cref{prop:pet-ch11-gh-distance-zero-isometric} the equivalence
  classes (up to isometry) form a genuine metric space $(\mathcal{M},\dGH)$.
\end{remark}

\begin{example}[$\varepsilon$-dense finite approximation]
  \label{ex:pet-ch11-epsilon-dense-approximation}
  \lean{PetersenLib.epsilonDense,PetersenLib.ghDistance_le_epsilon_of_epsilonDense}
  \source{petersen:page-0410,petersen:page-0411}
  \uses{def:pet-ch11-gromov-hausdorff-distance, def:pet-ch11-hausdorff-distance}
  \dcref{ch11:11.1.4}
  For compact $X$ and finite $A\subset X$ with $d_H(A,X)\le\varepsilon$
  (using the induced metric on $A$), call $A$ \emph{$\varepsilon$-dense} in
  $X$. Then $\dGH(X,A)\le\varepsilon$; since $X$ is compact, finite
  $\varepsilon$-dense subsets exist for every $\varepsilon>0$.
\end{example}
\begin{proof}
  \uses{ex:pet-ch11-epsilon-dense-approximation}
  \sketch
  For \(\delta>0\), define the cross-distance on \(X\sqcup A\) by
  \(|xa|=\delta+|xa|_X\). This admissible metric has
  \(d_H(X,A)<\varepsilon+\delta\); let \(\delta\downarrow0\).
\end{proof}


\begin{example}[correspondences bound $\dGH$ by $3\varepsilon$]
  \label{ex:pet-ch11-finite-approximation-3eps}
  \lean{PetersenLib.ghDistance_le_threeEpsilon_of_correspondence}
  \source{petersen:page-0411,petersen:page-0412}
  \uses{ex:pet-ch11-epsilon-dense-approximation}
  \dcref{ch11:11.1.5}
  Take finite sets
  $A=\{x_1,\dots,x_k\}\subset X$ and $B=\{y_1,\dots,y_k\}\subset Y$ that are
  $\varepsilon$-dense and satisfy
  $\big||x_ix_j|-|y_iy_j|\big|\le\varepsilon$ for $1\le i,j\le k$. Then
  $\dGH(X,Y)\le3\varepsilon$.
\end{example}
\begin{proof}
  \uses{ex:pet-ch11-finite-approximation-3eps}
  \sketch
  By \cref{ex:pet-ch11-epsilon-dense-approximation}, it remains to compare
  the finite approximations. On \(A\sqcup B\), set
  \[
    |x_i y_j|=\min_k\bigl(|x_i x_k|+\varepsilon+|y_jy_k|\bigr).
  \]
  The correspondence bound supplies the mixed triangle inequalities, so this is admissible and makes paired points at most \(\varepsilon\) apart. Thus \(\dGH(A,B)\le\varepsilon\). Since \(A\) and \(B\) are each \(\varepsilon\)-close to the ambient space, the triangle inequality gives \(\dGH(X,Y)\le3\varepsilon\).
\end{proof}


\begin{example}[lens spaces collapse to $S^2(1/2)$]
  \label{ex:pet-ch11-lens-spaces-limit}
  \lean{PetersenLib.lensSpace_ghConverges_toSphere}
  \source{petersen:page-0412}
  \uses{ex:pet-ch11-finite-approximation-3eps}
  \dcref{ch11:11.1.6}
  Equip $M_k=S^3/\mathbb{Z}_k$ with the metric induced by $S^3(1)$. The
  Riemannian submersion $M_k\to S^2(1/2)$ has fibres of diameter
  $2\pi/k\to0$; using \cref{ex:pet-ch11-finite-approximation-3eps} it follows
  that $M_k\to S^2(1/2)$ in the Gromov--Hausdorff topology.
\end{example}

\begin{example}[Berger spheres collapse to $S^2(1/2)$]
  \label{ex:pet-ch11-berger-spheres-limit}
  \lean{PetersenLib.bergerSphere_ghConverges_toSphere}
  \source{petersen:page-0412}
  \uses{ex:pet-ch11-lens-spaces-limit}
  \dcref{ch11:11.1.7}
  The Berger family also satisfies $(S^3,g_\varepsilon)\to S^2(1/2)$ as
  $\varepsilon\to0$. In both examples the volume goes to zero while the
  curvatures and diameters stay uniformly bounded; in the Berger case the
  manifolds are even simply connected, yet the topology changes drastically —
  in the lens-space case the fundamental groups of the sequence elements are
  even mutually different.
\end{example}

\begin{proposition}[$(\mathcal{M},\dGH)$ separable and complete]
  \label{prop:pet-ch11-gh-space-separable-complete}
  \lean{PetersenLib.gromovHausdorffSpace_separable_complete}
  \source{petersen:page-0412,petersen:page-0413}
  \uses{rem:pet-ch11-gh-pseudo-metric-quotient}
  \dcref{ch11:11.1.8}
  The quotient metric space $(\mathcal{M},\dGH)$ is both separable and
  complete.
\end{proposition}
\begin{proof}
  \uses{prop:pet-ch11-gh-space-separable-complete}
  \sketch
  Finite metric spaces are dense by
  \cref{ex:pet-ch11-epsilon-dense-approximation}, and perturbing their
  distance matrices to rational ones gives a countable dense family. For completeness, take a subsequence with
  \(\dGH(X_i,X_{i+1})<2^{-i}\), choose corresponding admissible metrics, and glue them along the chain. The completion of the resulting union contains compact closures \(Y_i\) isometric to \(X_i\), with \(d_H(Y_i,Y_{i+1})<2^{-i}\). These closures converge in Hausdorff distance to a compact set \(Y\), hence \(X_i\to Y\) in \(\dGH\).
\end{proof}


\subsection{Pointed Convergence}

\begin{definition}[pointed Gromov--Hausdorff distance]
  \label{def:pet-ch11-pointed-gh-distance}
  \lean{PetersenLib.pointedGromovHausdorffDistance}
  \source{petersen:page-0414}
  \uses{def:pet-ch11-gromov-hausdorff-distance, def:pet-ch11-hausdorff-distance}
  For pointed metric spaces $(X,x)$ and $(Y,y)$, define their \emph{pointed
  Gromov--Hausdorff distance} by
  \[
    \dGH\big((X,x),(Y,y)\big)=\inf\{d_H(X,Y)+|xy|\},
  \]
  the infimum over admissible metrics on $X\cup Y$ of the Hausdorff distance
  plus the distance between the marked points.
\end{definition}

\begin{definition}[pointed Gromov--Hausdorff topology]
  \label{def:pet-ch11-pointed-gh-topology}
  \lean{PetersenLib.PointedGromovHausdorffConvergence}
  \source{petersen:page-0414}
  \uses{def:pet-ch11-pointed-gh-distance}
  For $\mathcal{M}_*=\{(X,x,|\cdot|)\}$, the collection of proper pointed
  metric spaces, $(X_i,x_i,|\cdot|_i)\to(X,x,|\cdot|)$ in the \emph{pointed
  Gromov--Hausdorff topology} if for every $R$ there is a sequence $R_i\to R$
  such that the closed balls
  $\big(\overline{B}(x_i,R_i),x_i,|\cdot|_i\big)\to\big(\overline{B}(x,R),x,|\cdot|\big)$
  converge in the pointed Gromov--Hausdorff metric of
  \cref{def:pet-ch11-pointed-gh-distance}.
\end{definition}

\subsection{Convergence of Maps}

\begin{definition}[convergence of maps]
  \label{def:pet-ch11-convergence-of-maps}
  \lean{PetersenLib.MapsConvergeTo}
  \source{petersen:page-0414}
  \uses{def:pet-ch11-pointed-gh-topology}
  For maps $f_k:X_k\to Y_k$ with $X_k\to X$ and $Y_k\to Y$ in the
  Gromov--Hausdorff sense,
  $f_k$ \emph{converges} to $f:X\to Y$ if for every sequence $x_k\in X_k$
  with $x_k\to x\in X$ one has $f_k(x_k)\to f(x)$. This depends on the choice
  of Hausdorff-convergent ambient metrics. It is a strong notion: if
  $X_k=X$, $Y_k=Y$, $f_k=f$ for all $k$, then $f$ converges to itself only if
  $f$ is continuous.
\end{definition}

\begin{remark}[convergence of maps as an extension problem]
  \label{rem:pet-ch11-map-convergence-as-extension}
  \lean{PetersenLib.mapConvergence_iff_continuousExtension}
  \source{petersen:page-0415}
  \uses{def:pet-ch11-convergence-of-maps}
  This convergence preserves distance preservation and the submetry property. A
  sequence $f_k:X_k\to Y_k$ can be regarded as one continuous map
  $F:\bigcup_iX_i\to Y\cup\bigcup_iY_i$; the sequence converges iff $F$ has a
  continuous extension $X\cup\bigcup_iX_i\to Y\cup\bigcup_iY_i$ restricting to
  the limit map on $X$. When the $X_i$ are compact this is equivalent to $F$
  being uniformly continuous.
\end{remark}

\begin{definition}[equicontinuous family]
  \label{def:pet-ch11-equicontinuous-family}
  \lean{PetersenLib.EquicontinuousFamily}
  \source{petersen:page-0415}
  \uses{def:pet-ch11-convergence-of-maps}
  Call a family $f_k:X_k\to Y_k$ \emph{equicontinuous} when, for every
  $\varepsilon>0$ and $x_k\in X_k$ there is $\delta>0$ with
  $f_k(B(x_k,\delta))\subset B(f_k(x_k),\varepsilon)$ for all $k$. Uniformly
  Lipschitz families are equicontinuous.
\end{definition}

\begin{lemma}[Gromov--Hausdorff Arzelà--Ascoli]
  \label{lem:pet-ch11-arzela-ascoli-gh}
  \lean{PetersenLib.arzelaAscoli_gromovHausdorff}
  \source{petersen:page-0415}
  \uses{def:pet-ch11-equicontinuous-family}
  \dcref{ch11:11.1.9}
  Every equicontinuous family $f_k:X_k\to Y_k$ with $X_k\to X$ and $Y_k\to Y$
  in the (pointed) Gromov--Hausdorff topology, has a convergent subsequence.
  In the pointed case, one also assumes each $f_k$ preserves the base point.
\end{lemma}
\begin{proof}
  \uses{lem:pet-ch11-arzela-ascoli-gh}
  \sketch
  Choose dense sets \(A_i=\{a_j^i\}\subset X_i\) with \(a_j^i\to a_j\) and \(\{a_j\}\) dense in \(X\). Compactness of \(Y\) and a diagonal subsequence give limits of \(f_i(a_j^i)\) for every \(j\). Equicontinuity extends these values uniquely to a continuous \(f:X\to Y\) and yields convergence in the sense of \cref{def:pet-ch11-convergence-of-maps}.
\end{proof}


\subsection{Compactness of Classes of Metric Spaces}

\begin{definition}[capacity and covering functions]
  \label{def:pet-ch11-capacity-covering-functions}
  \lean{PetersenLib.capacityFunction,PetersenLib.coveringFunction}
  \source{petersen:page-0416}
  For compact $X$, write
  $\Cap(\varepsilon)=\Cap_X(\varepsilon)$ is the maximum number of disjoint
  $\varepsilon/2$-balls in $X$, and $\Cov(\varepsilon)=\Cov_X(\varepsilon)$
  the minimum number of $\varepsilon$-balls needed to cover $X$.
\end{definition}

\begin{lemma}[covering is bounded by capacity]
  \label{lem:pet-ch11-cov-le-cap}
  \lean{PetersenLib.coveringFunction_le_capacityFunction}
  \source{petersen:page-0416}
  \uses{def:pet-ch11-capacity-covering-functions}
  The covering and capacity numbers satisfy
  $\Cov(\varepsilon)\le\Cap(\varepsilon)$.
\end{lemma}
\begin{proof}
  \uses{lem:pet-ch11-cov-le-cap}
  \sketch
  Take a maximal disjoint family of \(\varepsilon/2\)-balls. Maximality forces
  the concentric \(\varepsilon\)-balls to cover \(X\), so their number gives
  \(\Cov_X(\varepsilon)\le\Cap_X(\varepsilon)\).
\end{proof}


\begin{lemma}[covering/capacity under Gromov--Hausdorff proximity]
  \label{lem:pet-ch11-gh-close-cov-cap-comparison}
  \lean{PetersenLib.coveringFunction_capacityFunction_ghClose}
  \source{petersen:page-0416}
  \uses{def:pet-ch11-capacity-covering-functions, def:pet-ch11-gromov-hausdorff-distance}
  Whenever $\dGH(X,Y)<\delta$, one has
  $\Cov_X(\varepsilon+2\delta)\le\Cov_Y(\varepsilon)$ and
  $\Cap_X(\varepsilon)\ge\Cap_Y(\varepsilon+2\delta)$.
\end{lemma}
\begin{proof}
  \uses{lem:pet-ch11-gh-close-cov-cap-comparison}
  \sketch
  Realize \(d_H(X,Y)<\delta\) in an admissible metric. Replacing each center
  of an \(\varepsilon\)-cover of \(Y\) by a point of \(X\) within \(\delta\)
  gives an \((\varepsilon+2\delta)\)-cover of \(X\). Likewise, centers of
  disjoint \((\varepsilon+2\delta)/2\)-balls in \(Y\) map to centers of
  disjoint \(\varepsilon/2\)-balls in \(X\), proving the capacity inequality.
\end{proof}


\begin{proposition}[precompactness criterion]
  \label{prop:pet-ch11-gromov-precompactness-criterion}
  \lean{PetersenLib.gromovPrecompactnessCriterion}
  \source{petersen:page-0416,petersen:page-0417}
  \uses{lem:pet-ch11-cov-le-cap, lem:pet-ch11-gh-close-cov-cap-comparison}
  \dcref{ch11:11.1.10}
  Let $\mathcal{C}\subset(\mathcal{M},\dGH)$ have diameters bounded by
  $D<\infty$. The following are equivalent: (1) $\mathcal{C}$ is precompact;
  (2) there is
  $N_1:(0,\infty)\to(0,\infty)$ with $\Cap_X(\varepsilon)\le N_1(\varepsilon)$
  for all $X\in\mathcal{C}$; (3) there is $N_2:(0,\infty)\to(0,\infty)$ with
  $\Cov_X(\varepsilon)\le N_2(\varepsilon)$ for all $X\in\mathcal{C}$.
\end{proposition}
\begin{proof}
  \uses{prop:pet-ch11-gromov-precompactness-criterion}
  \sketch
  A finite \(\varepsilon/4\)-net for a precompact class and
  \cref{lem:pet-ch11-gh-close-cov-cap-comparison} give uniform capacity
  bounds; \cref{lem:pet-ch11-cov-le-cap} then gives uniform covering bounds.
  Conversely, approximate each \(X\) by an \(\varepsilon/2\)-dense finite set
  of uniformly bounded cardinality. Its distance matrix lies in a bounded
  finite-dimensional cube, which has a finite \(\varepsilon/2\)-net; the
  correspondence estimate supplies a finite \(3\varepsilon\)-net for the class.
\end{proof}


\begin{corollary}[pointed precompactness]
  \label{cor:pet-ch11-pointed-precompactness-criterion}
  \lean{PetersenLib.pointedPrecompactnessCriterion}
  \source{petersen:page-0417}
  \uses{prop:pet-ch11-gromov-precompactness-criterion, def:pet-ch11-pointed-gh-topology}
  \dcref{ch11:11.1.11}
  A class $\mathcal{C}\subset\mathcal{M}_*$ is precompact exactly when, for
  every $R>0$, the
  collection $\{\overline{B}(x,R)\mid(X,x)\in\mathcal{C}\}\subset(\mathcal{M},\dGH)$
  is precompact.
\end{corollary}

\begin{definition}[metric doubling condition]
  \label{def:pet-ch11-metric-doubling-condition}
  \lean{PetersenLib.MetricDoublingCondition}
  \source{petersen:page-0417}
  Say that $X$ has the \emph{metric doubling condition} with constant $C$ when
  every ball $B(p,R)$ can be covered by at most $C$ balls of radius $R/2$.
\end{definition}

\begin{proposition}[doubling implies precompact]
  \label{prop:pet-ch11-doubling-implies-precompact}
  \lean{PetersenLib.doublingCondition_implies_precompact}
  \source{petersen:page-0417}
  \uses{def:pet-ch11-metric-doubling-condition, prop:pet-ch11-gromov-precompactness-criterion}
  \dcref{ch11:11.1.12}
  A family $\mathcal{C}\subset(\mathcal{M},\dGH)$ with a common doubling
  constant $C<\infty$ and $\diam X\le D<\infty$ for every member, then
  $\mathcal{C}$ is precompact.
\end{proposition}
\begin{proof}
  \uses{prop:pet-ch11-doubling-implies-precompact}
  \sketch
  Iterating the doubling hypothesis covers every \(X\) by at most \(C^N\)
  balls of radius \(2^{-N}D\). Choosing \(N\) from \(\varepsilon\) yields a
  uniform covering function, so
  \cref{prop:pet-ch11-gromov-precompactness-criterion} applies.
\end{proof}


\begin{corollary}[Ricci/diameter precompactness]
  \label{cor:pet-ch11-ricci-diam-precompactness}
  \lean{PetersenLib.riccciBounded_precompact}
  \source{petersen:page-0417}
  \uses{prop:pet-ch11-doubling-implies-precompact, cor:pet-ch11-pointed-precompactness-criterion}
  \dcref{ch11:11.1.13}
  For $n\ge2$, $k\in\mathbb{R}$, and $D>0$, the following classes are
  precompact: (1) closed Riemannian
  $n$-manifolds with $\Ric\ge(n-1)k$, $\diam\le D$; (2) the
  class of pointed complete Riemannian $n$-manifolds with $\Ric\ge(n-1)k$ is
  precompact.
\end{corollary}
\begin{proof}
  \uses{cor:pet-ch11-ricci-diam-precompactness}
  \sketch
  If \(N\) disjoint \(\varepsilon\)-balls lie in \(\overline B(x,R)\), choose one of minimal volume. Relative volume comparison gives
  \[
    N\le \frac{\vol B(x,R)}{\vol B(x_i,\varepsilon)}
      \le \frac{v(n,k,2R)}{v(n,k,\varepsilon)}.
  \]
  This uniform local capacity bound and
  \cref{cor:pet-ch11-pointed-precompactness-criterion} prove pointed
  precompactness. Equivalently, one may use the doubling control in
  \cref{prop:pet-ch11-doubling-implies-precompact}; a diameter bound makes the
  convergence global.
\end{proof}


\section{Hölder Spaces and Schauder Estimates}

\subsection{Hölder Spaces}

\begin{definition}[$C^0$ and $C^m$ norms]
  \label{def:pet-ch11-cm-norm}
  \lean{PetersenLib.cZeroNorm,PetersenLib.cmNorm}
  \source{petersen:page-0418,petersen:page-0419}
  On a bounded domain $\Omega\subset\mathbb{R}^n$, let
  $C^0(\Omega,\mathbb{R}^k)$ be the bounded continuous maps with sup-norm
  $\|u\|_{C^0}=\sup_{x\in\Omega}|u(x)|$, a Banach space. Using multi-index
  notation $\partial^Iu=\partial_1^{i_1}\cdots\partial_n^{i_n}u$ for
  $I=(i_1,\dots,i_n)$, $|I|=i_1+\dots+i_n$, define $C^m(\Omega,\mathbb{R}^k)$ as
  the functions with $m$ continuous partial derivatives, with norm
  \[
    \|u\|_{C^m}=\|u\|_{C^0}+\sum_{1\le|I|\le m}\|\partial^Iu\|_{C^0}.
  \]
  This makes $C^m(\Omega,\mathbb{R}^k)$ a Banach space, but the inclusions
  $C^m\subset C^{m-1}$ are not closed subspaces (e.g.\ $f(x)=|x|$ lies in the
  $C^0([-1,1],\mathbb{R})$-closure of $C^1([-1,1],\mathbb{R})$).
\end{definition}

\begin{definition}[Hölder pseudo-norm]
  \label{def:pet-ch11-holder-pseudonorm}
  \lean{PetersenLib.holderPseudonorm}
  \source{petersen:page-0419}
  \uses{def:pet-ch11-cm-norm}
  Given $\alpha\in(0,1]$ and $u:\Omega\to\mathbb{R}^k$, set the
  \emph{$C^\alpha$ pseudo-norm}
  given by
  \[
    \|u\|_\alpha=\sup_{x,y\in\Omega}\frac{|u(x)-u(y)|}{|x-y|^\alpha}.
  \]
  When $\alpha=1$ this is the best Lipschitz constant for $u$.
\end{definition}

\begin{definition}[Hölder space $C^{m,\alpha}$]
  \label{def:pet-ch11-holder-space}
  \lean{PetersenLib.holderSpace}
  \source{petersen:page-0419}
  \uses{def:pet-ch11-cm-norm, def:pet-ch11-holder-pseudonorm}
  Define the \emph{Hölder space} $C^{m,\alpha}(\Omega,\mathbb{R}^k)$ as those
  $u\in C^m(\Omega,\mathbb{R}^k)$ whose $m$th-order partial derivatives all have
  finite $C^\alpha$ pseudo-norm, with norm
  \[
    \|u\|_{C^{m,\alpha}}=\|u\|_{C^m}+\sum_{|I|=m}\|\partial^Iu\|_\alpha.
  \]
  If the domain needs to be specified, one writes $\|u\|_{C^{m,\alpha},\Omega}$.
\end{definition}

\begin{lemma}[Hölder spaces are Banach, compactly nested]
  \label{lem:pet-ch11-holder-space-banach-compact-inclusion}
  \lean{PetersenLib.holderSpace_banach,PetersenLib.holderSpace_inclusion_compact}
  \source{petersen:page-0419,petersen:page-0420}
  \uses{def:pet-ch11-holder-space}
  \dcref{ch11:11.2.1}
  The space $C^{m,\alpha}(\Omega,\mathbb{R}^k)$ is Banach, and for
  $\beta<\alpha$ its inclusion
  $C^{m,\alpha}(\Omega,\mathbb{R}^k)\subset C^{m,\beta}(\Omega,\mathbb{R}^k)$ is
  always compact (maps closed bounded sets to compact sets).
\end{lemma}
\begin{proof}
  \uses{lem:pet-ch11-holder-space-banach-compact-inclusion}
  \sketch
  It suffices to take \(m=0\). A \(C^\alpha\)-Cauchy sequence converges uniformly; passing to the limit in the Hölder quotients proves convergence in \(C^\alpha\). A bounded \(C^\alpha\) set is equicontinuous, hence relatively compact in \(C^0\) by Arzelà--Ascoli. The interpolation inequality
  \[
    [u]_\beta\le [u]_\alpha^{\beta/\alpha}(2\|u\|_{C^0})^{1-\beta/\alpha}
  \]
  upgrades \(C^0\)-convergence to \(C^\beta\)-convergence for \(\beta<\alpha\); apply this to derivatives for general \(m\).
\end{proof}


\subsection{Elliptic Estimates}

\begin{definition}[elliptic operator]
  \label{def:pet-ch11-elliptic-operator}
  \lean{PetersenLib.EllipticOperator}
  \source{petersen:page-0420,petersen:page-0421}
  \uses{def:pet-ch11-holder-pseudonorm}
  For $Lu=a^{ij}\partial_i\partial_ju+b^i\partial_iu$ with $a^{ij}=a^{ji}$,
  call $L$ \emph{elliptic} when $(a^{ij})$ is positive definite. Assume
  eigenvalues of $(a^{ij})$ lie in $[\lambda,\lambda^{-1}]$, $\lambda>0$, and
  $\|a^{ij}\|_\alpha\le\lambda^{-1}$, $\|b^i\|_\alpha\le\lambda^{-1}$.
\end{definition}

\begin{theorem}[Schauder estimates]
  \label{thm:pet-ch11-schauder-estimates}
  \lean{PetersenLib.schauderEstimates}
  \source{petersen:page-0421}
  \uses{def:pet-ch11-elliptic-operator, def:pet-ch11-holder-space}
  \dcref{ch11:11.2.2}
  Let $\Omega\subset\mathbb{R}^n$ have diameter $\le D$, and let
  $K\subset\Omega$ satisfy $d(K,\partial\Omega)\ge\delta$. For
  $\alpha\in(0,1)$ there is
  $C=C(n,\alpha,\lambda,\delta,D)$ with
  \[
    \|u\|_{C^{2,\alpha},K}\le C\big(\|Lu\|_{C^0,\Omega}+\|u\|_{C^0,\Omega}\big),
    \qquad
    \|u\|_{C^{1,\alpha},K}\le C\big(\|Lu\|_{C^0,\Omega}+\|u\|_{C^0,\Omega}\big).
  \]
  If in addition $\Omega$ has smooth boundary and $u=\varphi$ on
  $\partial\Omega$, there is $C=C(n,\alpha,\lambda,D)$ with
  \[
    \|u\|_{C^{2,\alpha},\Omega}\le C\big(\|Lu\|_{C^0,\Omega}+\|\varphi\|_{C^{2,\alpha},\partial\Omega}\big).
  \]
\end{theorem}

\begin{remark}[divergence form and the Laplace--Beltrami operator]
  \label{rem:pet-ch11-divergence-form-laplace-beltrami}
  \lean{PetersenLib.divergenceFormOperator,PetersenLib.laplaceBeltrami}
  \source{petersen:page-0421,petersen:page-0422}
  \uses{def:pet-ch11-elliptic-operator}
  The estimates in \cref{thm:pet-ch11-schauder-estimates} start with
  $Lu=\Delta u=\delta^{ij}\partial_i\partial_ju$,
  then, via a linear change of coordinates, for constant-coefficient
  operators $Lu=a^{ij}\partial_i\partial_ju$; Schauder's trick is that the
  hypotheses on $a^{ij}$ make it locally almost constant, and a partition of
  unity finishes the argument. The first-order term causes no extra trouble,
  especially when $L$ is in \emph{divergence form}
  $Lu=a^{ij}\partial_i\partial_ju+b^i\partial_iu=\partial_i(a^{ij}\partial_ju)$,
  which integrates by parts as
  $\int_\Omega\partial_i(a^{ij}\partial_ju)\,h=-\int_\Omega a^{ij}\partial_iu\,\partial_ih$
  for $h=0$ on $\partial\Omega$. The Laplacian of a Riemannian metric $g$ is of
  this form in local coordinates:
  \[
    Lu=\Delta_gu=\frac1{\sqrt{\det g_{ij}}}\partial_i\big(\sqrt{\det g_{ij}}\,g^{ij}\partial_ju\big),
  \]
  so $\int v\,Lu\,\vol=\int v\,\partial_i\big(\sqrt{\det g_{ij}}\,g^{ij}\partial_ju\big)$.
\end{remark}

\begin{corollary}[higher-order Schauder estimates]
  \label{cor:pet-ch11-higher-schauder-estimates}
  \lean{PetersenLib.schauderEstimates_higherOrder}
  \source{petersen:page-0422}
  \uses{thm:pet-ch11-schauder-estimates}
  \dcref{ch11:11.2.3}
  When $\|a^{ij}\|_{C^{m,\alpha}},\|b^i\|_{C^{m,\alpha}}\le\lambda^{-1}$,
  one has $C=C(n,m,\alpha,\lambda,\delta,D)$ with
  \[
    \|u\|_{C^{m+2,\alpha},K}\le C\big(\|Lu\|_{C^{m,\alpha},\Omega}+\|u\|_{C^0,\Omega}\big),
  \]
  and on a domain with smooth boundary,
  $\|u\|_{C^{m+2,\alpha},\Omega}\le C\big(\|Lu\|_{C^{m,\alpha},\Omega}+\|\varphi\|_{C^{m+2,\alpha},\partial\Omega}\big)$.
\end{corollary}

\begin{theorem}[Dirichlet problem]
  \label{thm:pet-ch11-dirichlet-problem-unique-solution}
  \lean{PetersenLib.dirichletProblem_existsUnique}
  \source{petersen:page-0422}
  \uses{thm:pet-ch11-schauder-estimates}
  \dcref{ch11:11.2.4}
  On a bounded smooth domain $\Omega\subset\mathbb{R}^n$, the problem
  $Lu=f$, $u|_{\partial\Omega}=\varphi$ has a unique
  solution $u\in C^{2,\alpha}(\Omega)$ when $f\in C^\alpha(\Omega)$,
  $\varphi\in C^{2,\alpha}(\partial\Omega)$.
\end{theorem}
\begin{proof}
  \uses{thm:pet-ch11-dirichlet-problem-unique-solution}
  \sketch
  The Schauder theory supplies a solution. If \(u_1,u_2\) solve the same data, \(w=u_1-u_2\) satisfies \(Lw=0\) and vanishes on \(\partial\Omega\); the maximum principle gives \(w=0\).
\end{proof}


\subsection{Harmonic Coordinates}

\begin{lemma}[existence of harmonic coordinates]
  \label{lem:pet-ch11-harmonic-coordinates-exist}
  \lean{PetersenLib.harmonicCoordinates_exist}
  \source{petersen:page-0422,petersen:page-0423}
  \uses{thm:pet-ch11-schauder-estimates, rem:pet-ch11-divergence-form-laplace-beltrami}
  \dcref{ch11:11.2.5}
  Every point $p$ of an $n$-dimensional Riemannian manifold has a neighborhood
  $U$ with a \emph{harmonic coordinate system}
  $x=(x^1,\dots,x^n):U\to\mathbb{R}^n$, i.e.\ each $x^i$ is harmonic for
  $\Delta_g$.
\end{lemma}
\begin{proof}
  \uses{lem:pet-ch11-harmonic-coordinates-exist}
  \sketch
  Start with normal coordinates \(y\) on a small Euclidean ball and solve
  \(\Delta_gx^k=0\) with \(x^k=y^k\) on the boundary by
  \cref{thm:pet-ch11-dirichlet-problem-unique-solution}. The divergence form
  in \cref{rem:pet-ch11-divergence-form-laplace-beltrami} and the estimates
  of \cref{thm:pet-ch11-schauder-estimates} give
  \[
    \|x-y\|_{C^{2,\alpha}}\le C\|\Delta_gy\|_{C^\alpha}.
  \]
  In normal coordinates \(\Delta_gy(0)=0\), so the right side tends to zero with the radius. Thus \(Dx\) remains invertible and \(x\) is a harmonic coordinate system.
\end{proof}


\begin{lemma}[Laplacian and Ricci curvature in harmonic coordinates]
  \label{lem:pet-ch11-laplacian-ricci-harmonic-coords}
  \lean{PetersenLib.laplacian_eq_gijPartial_harmonic,PetersenLib.ricci_harmonicCoords_formula}
  \source{petersen:page-0423,petersen:page-0424,petersen:page-0425}
  \uses{lem:pet-ch11-harmonic-coordinates-exist}
  \dcref{ch11:11.2.6}
  In harmonic coordinates $x:U\to\mathbb{R}^n$ for $(M,g)$,
  \begin{equation}
    \Delta u=\frac1{\sqrt{\det g_{ij}}}\partial_i\big(\sqrt{\det g_{ij}}\,g^{ij}\partial_ju\big)=g^{ij}\partial_i\partial_ju,
  \end{equation}
  \begin{equation}
    \tfrac12\Delta g_{ij}+Q(g,\partial g)=-\Ric_{ij}=-\Ric(\partial_i,\partial_j),
  \end{equation}
  where $Q$ is a universal rational expression, polynomial in $g$ and
  quadratic in $\partial g$ in the numerator, with denominator a power of
  $\sqrt{\det g_{ij}}$.
\end{lemma}
\begin{proof}
  \uses{lem:pet-ch11-laplacian-ricci-harmonic-coords}
  \sketch
  Harmonicity of \(x^k\) makes
  \(\partial_i(\sqrt{\det g}\,g^{ik})=0\), so the divergence formula reduces to
  \(\Delta u=g^{ij}\partial_i\partial_ju\). Substituting the Christoffel symbols into the Ricci formula and using these identities gives
  \[
    \tfrac12\Delta g_{ij}+Q(g,\partial g)=-\Ric_{ij},
  \]
  where \(Q\) has the stated universal rational form. For a harmonic function, the Bochner identity
  \(\Delta(\tfrac12|\nabla u|^2)=|\Hess u|^2+\Ric(\nabla u,\nabla u)\)
  supplies the auxiliary estimates used in the coordinate calculation.
\end{proof}


\begin{remark}[Einstein metrics in harmonic coordinates are smooth]
  \label{rem:pet-ch11-einstein-harmonic-regularity-bootstrap}
  \lean{PetersenLib.einsteinHarmonicWeakSolution_smooth}
  \source{petersen:page-0425,petersen:page-0426}
  \uses{lem:pet-ch11-laplacian-ricci-harmonic-coords, cor:pet-ch11-higher-schauder-estimates}
  For an Einstein metric $\Ric_{ij}=(n-1)kg_{ij}$,
  \cref{lem:pet-ch11-laplacian-ricci-harmonic-coords} reads
  $\tfrac12\Delta g_{ij}=-(n-1)kg_{ij}-Q(g,\partial g)$, which makes sense
  once $g_{ij}\in C^1$ and, after multiplying by a test function and
  integrating by parts, can be interpreted weakly. If $g_{ij}\in C^{1,\alpha}$
  the right side lies in some $C^\beta$, so \cref{cor:pet-ch11-higher-schauder-estimates}
  gives $g_{ij}\in C^{2,\beta}$; bootstrapping shows $g_{ij}\in C^\infty$
  (in fact real-analytic). Hence any metric that is, in harmonic coordinates,
  a weak solution of the Einstein equation is automatically smooth.
\end{remark}

\section{Norms and Convergence of Manifolds}

\subsection{Norms of Riemannian Manifolds}

\begin{definition}[$C^{m,\alpha}$ norm at scale $r$]
  \label{def:pet-ch11-cmalpha-norm-pointed}
  \lean{PetersenLib.pointedNorm}
  \source{petersen:page-0426,petersen:page-0427}
  \uses{def:pet-ch11-holder-space}
  Define $\|(M,g,p)\|_{C^{m,\alpha},r}\le Q$ when a $C^{m+1,\alpha}$ chart
  exists from $\varphi:(B(0,r),0)\subset\mathbb{R}^n\to(U,p)\subset M$ such
  that
  \begin{equation}
    \begin{aligned}
      \text{(n1)}\quad
      &|D\varphi|\le e^Q\text{ on }B(0,r),\\
      &|D\varphi^{-1}|\le e^Q\text{ on }U,\\
      &e^{-2Q}\delta_{kl}v^kv^l\le g_{kl}v^kv^l,\\
      &g_{kl}v^kv^l\le e^{2Q}\delta_{kl}v^kv^l,\\
      &v\in\mathbb{R}^n.
    \end{aligned}
  \end{equation}
  \begin{equation}
    \begin{aligned}
      \text{(n2)}\quad
      &r^{|I|+\alpha}\|\partial^Ig_{ij}\|_\alpha\le Q\\
      &\text{for every multi-index }I,\\
      &0\le|I|\le m.
    \end{aligned}
  \end{equation}
\end{definition}

\begin{definition}[global $C^{m,\alpha}$ norm of a manifold]
  \label{def:pet-ch11-cmalpha-norm-global}
  \lean{PetersenLib.globalNorm}
  \source{petersen:page-0427}
  \uses{def:pet-ch11-cmalpha-norm-pointed}
  $\|(M,g)\|_{C^{m,\alpha},r}=\sup_{p\in M}\|(M,g,p)\|_{C^{m,\alpha},r}$. The
  charts are regarded as maps from the fixed model $B(0,r)$ into $M$, so that
  the norm measures how far the manifold differs from a Euclidean ball: (n1)
  bounds the metric coefficients from above and below in the chosen
  coordinates (giving uniform ellipticity for the Laplacian), and (n2)
  bounds the derivatives of the metric. The definition makes sense for
  merely $C^{m,\alpha}$ metrics, without assuming more smoothness; some
  constructions (e.g.\ exponential maps) may then be ill-defined if $m\le1$.
  The pointed norm is always finite, but the global norm need not be finite
  on any scale when $M$ is noncompact.
\end{definition}

\begin{example}[flat manifolds have norm zero]
  \label{ex:pet-ch11-flat-manifold-norm-zero}
  \lean{PetersenLib.flatManifold_norm_zero}
  \source{petersen:page-0427}
  \uses{def:pet-ch11-cmalpha-norm-global}
  \dcref{ch11:11.3.1}
  A complete flat manifold satisfies $\|(M,g)\|_{C^{m,\alpha},r}=0$
  for all $r\le\inj(M,g)$; in particular
  $\|(\mathbb{R}^n,g_{\mathrm{eu}})\|_{C^{m,\alpha},r}=0$ for all $r$.
\end{example}

\subsection{Convergence of Riemannian Manifolds}

\begin{definition}[$C^{m,\alpha}$ convergence of tensors on a compact subset]
  \label{def:pet-ch11-tensor-convergence-compact-subset}
  \lean{PetersenLib.tensorConverges_onCompactSubset}
  \source{petersen:page-0427}
  On a fixed closed manifold $M$ (or precompact $A\subset M$), a sequence
  of functions on $A$ \emph{converges in $C^{m,\alpha}$} if it converges in
  the charts of some fixed finite covering of coordinate patches that are
  uniformly bi-Lipschitz; this is independent of the covering. A sequence of
  tensors converges in $C^{m,\alpha}$ if its components converge in these
  patches, giving a notion of convergence of Riemannian metrics on compact
  subsets of a fixed manifold.
\end{definition}

\begin{definition}[pointed $C^{m,\alpha}$ convergence of manifolds]
  \label{def:pet-ch11-pointed-cmalpha-convergence}
  \lean{PetersenLib.PointedCmAlphaConvergence}
  \source{petersen:page-0427,petersen:page-0428}
  \uses{def:pet-ch11-tensor-convergence-compact-subset, def:pet-ch11-cmalpha-norm-pointed}
  We say pointed complete Riemannian manifolds converge in the
  \emph{pointed $C^{m,\alpha}$ topology}, $(M_i,g_i,p_i)\to(M,g,p)$, if for
  every $R>0$ there is a domain $\Omega\supset B(p,R)\subset M$ and
  embeddings $F_i:\Omega\to M_i$ for large $i$ with $F_i(p)=p_i$,
  $F_i(\Omega)\supset B(p_i,R)$, and $F_i^*g_i\to g$ on $\Omega$ in the
  $C^{m,\alpha}$ topology.
\end{definition}

\begin{remark}[pointed $C^{m,\alpha}$ convergence implies GH convergence]
  \label{rem:pet-ch11-cmalpha-convergence-implies-gh}
  \lean{PetersenLib.cmAlphaConvergence_implies_ghConvergence}
  \source{petersen:page-0428}
  \uses{def:pet-ch11-pointed-cmalpha-convergence, def:pet-ch11-pointed-gh-topology}
  Pointed $C^{m,\alpha}$ convergence entails pointed Gromov--Hausdorff
  convergence. When all manifolds are closed with uniformly bounded
  diameter, the $F_i$ are diffeomorphisms, so one may speak of unpointed
  convergence, which forces all manifolds in the tail of the sequence to be
  diffeomorphic; classes of closed manifolds precompact in some $C^{m,\alpha}$
  topology contain at most finitely many diffeomorphism types.
\end{remark}

\begin{remark}[pullback convergence versus classical convergence]
  \label{rem:pet-ch11-cmalpha-convergence-pullback-warning}
  \lean{PetersenLib.pullbackConvergence_not_classicalConvergence}
  \source{petersen:page-0428}
  \uses{def:pet-ch11-pointed-cmalpha-convergence}
  Pullbacks $g_i=F_i^*g$ by diffeomorphisms $F_i$ of a fixed $M$ may converge
  in the pointed $C^{m,\alpha}$ sense of
  \cref{def:pet-ch11-pointed-cmalpha-convergence} without converging in any
  fixed coordinate system: e.g., on $M=S^2$ with the standard metric $g$, let
  $F_i$ be diffeomorphisms from the flow of a vector field vanishing only at
  the poles and pointing toward the south pole. Away from the poles,
  $F_i^*g\to0$ in fixed coordinates as the flow collapses the sphere toward
  the south pole, so $(S^2,F_i^*g)$ does not converge classically; but
  $(S^2,F_i^*g)$ is isometric to $(S^2,g)$ for every $i$, so as pointed
  Riemannian manifolds up to isometry the sequence is constant.
\end{remark}

\subsection{Properties of the Norm}

\begin{proposition}[elementary properties of the norm]
  \label{prop:pet-ch11-norm-properties}
  \lean{PetersenLib.norm_scaling,PetersenLib.norm_monotone_continuous,PetersenLib.norm_continuous_underConvergence,PetersenLib.norm_ballContainment,PetersenLib.norm_chart_attained,PetersenLib.norm_sup_attained}
  \source{petersen:page-0428,petersen:page-0429,petersen:page-0430,petersen:page-0431}
  \uses{def:pet-ch11-cmalpha-norm-pointed, def:pet-ch11-cmalpha-norm-global, def:pet-ch11-pointed-cmalpha-convergence}
  \dcref{ch11:11.3.2}
  For $(M,g,p)$, $m\ge0$, and $\alpha\in(0,1]$, the norm has:
  \begin{enumerate}
  \item[(1)] $\|(M,g,p)\|_{C^{m,\alpha},r}=\|(M,\lambda^2g,p)\|_{C^{m,\alpha},\lambda r}$ for all $\lambda>0$.
  \item[(2)] $r\mapsto\|(M,g,p)\|_{C^{m,\alpha},r}$ is increasing, continuous, and $\to0$ as $r\to0$.
  \item[(3)] If $(M_i,g_i,p_i)\to(M,g,p)$ in $C^{m,\alpha}$, then
    $\|(M_i,g_i,p_i)\|_{C^{m,\alpha},r}\to\|(M,g,p)\|_{C^{m,\alpha},r}$ for all $r>0$;
    if in addition all manifolds have uniformly bounded diameter,
    $\|(M_i,g_i)\|_{C^{m,\alpha},r}\to\|(M,g)\|_{C^{m,\alpha},r}$ for all $r>0$.
  \item[(4)] If $\|(M,g,p)\|_{C^{m,\alpha},r}<Q$, then for $x_1,x_2\in B(0,r)$,
    $e^{-Q}\min\{|x_1-x_2|,2r-|x_1|-|x_2|\}\le|\varphi(x_1)\varphi(x_2)|\le e^Q|x_1-x_2|$.
  \item[(5)] The infimum $Q=\|(M,g,p)\|_{C^{m,\alpha},r}$ is attained by some
    chart satisfying (n1), (n2) with this $Q$.
  \item[(6)] When $M$ is compact, the supremum $\|(M,g)\|_{C^{m,\alpha},r}=\sup_p\|(M,g,p)\|_{C^{m,\alpha},r}$ is a maximum.
  \end{enumerate}
\end{proposition}
\begin{proof}
  \uses{prop:pet-ch11-norm-properties}
  \sketch
  Scaling a chart proves (1). Restriction to smaller balls gives monotonicity; rescaling the domain and comparing the defining bounds gives continuity and the limit \(0\) at the origin. Pullback along the embeddings in pointed convergence yields both semicontinuity inequalities in (3), and bounded diameter upgrades them to the global norm. The derivative bounds in (n1) give the upper distance estimate; a minimizing segment either remains in the chart or meets its boundary, giving the lower estimate in (4). Arzelà--Ascoli applied to minimizing charts proves attainment in (5), and continuity in the base point plus compactness proves (6).
\end{proof}


\begin{corollary}[ball containment in the chart domain]
  \label{cor:pet-ch11-norm-bound-ball-containment}
  \lean{PetersenLib.norm_ball_subset_chartDomain}
  \source{petersen:page-0431}
  \uses{prop:pet-ch11-norm-properties}
  \dcref{ch11:11.3.3}
  If $\|(M,g,p)\|_{C^{m,\alpha},r}\le Q$, then $B(p,e^{-Q}r)\subset U$.
\end{corollary}
\begin{proof}
  \uses{cor:pet-ch11-norm-bound-ball-containment}
  \sketch
  Let \(q\in\partial U\) minimize distance to \(p\). A minimizing segment from \(p\) to \(q\) has coordinate radius tending to \(r\); the lower distance bound in part (4) of \cref{prop:pet-ch11-norm-properties} therefore gives
  \[
    |pq|\ge e^{-Q}r.
  \]
  Hence \(B(p,e^{-Q}r)\subset U\).
\end{proof}


\begin{corollary}[norm zero implies flat]
  \label{cor:pet-ch11-norm-zero-implies-flat}
  \lean{PetersenLib.norm_zero_implies_flat}
  \source{petersen:page-0431}
  \uses{prop:pet-ch11-norm-properties, cor:pet-ch11-norm-bound-ball-containment}
  \dcref{ch11:11.3.4}
  If $\|(M,g,p)\|_{C^{m,\alpha},r}=0$ for some $r$, then $p$ has a flat
  neighborhood.
\end{corollary}
\begin{proof}
  \uses{cor:pet-ch11-norm-zero-implies-flat}
  \sketch
  Norm zero gives an attained chart with \(Q=0\) and, by the preceding corollary, a neighborhood containing \(B(p,r)\). Its metric coefficients equal \(\delta_{ij}\), so the chart is a local Euclidean isometry and the neighborhood is flat.
\end{proof}


\subsection{The Harmonic Norm}

\begin{definition}[harmonic norm]
  \label{def:pet-ch11-harmonic-norm}
  \lean{PetersenLib.harmonicNorm}
  \source{petersen:page-0431,petersen:page-0432}
  \uses{def:pet-ch11-cmalpha-norm-pointed, lem:pet-ch11-harmonic-coordinates-exist}
  The \emph{harmonic norm} $\|(M,g,p)\|^{\mathrm{har}}_{C^{m,\alpha},r}$ is
  defined exactly as in \cref{def:pet-ch11-cmalpha-norm-pointed}, except that
  $\varphi^{-1}:U\to\mathbb{R}^n$ is additionally required to be harmonic for
  $g$, i.e.\ for each $j$,
  $\frac1{\sqrt{\det[g_{ij}]}}\partial_j\big(\sqrt{\det[g_{ij}]}\,g^{ij}\big)=0$.
\end{definition}

\begin{proposition}[harmonic-norm analogue]
  \label{prop:pet-ch11-harmonic-norm-properties}
  \lean{PetersenLib.harmonicNorm_properties}
  \source{petersen:page-0432}
  \uses{def:pet-ch11-harmonic-norm, prop:pet-ch11-norm-properties, lem:pet-ch11-harmonic-coordinates-exist}
  \dcref{ch11:11.3.5}
  \cref{prop:pet-ch11-norm-properties} holds for the harmonic norm when
  $m\ge1$.
\end{proposition}
\begin{proof}
  \uses{prop:pet-ch11-harmonic-norm-properties}
  \sketch
  Repeat \cref{prop:pet-ch11-norm-properties}. The harmonic coordinates from
  \cref{lem:pet-ch11-harmonic-coordinates-exist}, normalized by
  \(g_{ij}(p)=\delta_{ij}\), \(\partial_kg_{ij}(p)=0\), give the small-scale
  limit. Under convergence, harmonic charts converge to harmonic charts,
  giving one semicontinuity inequality. For the reverse inequality, solve the
  harmonic Dirichlet problem on a slightly larger limiting chart;
  \cref{cor:pet-ch11-higher-schauder-estimates} makes the solutions converge
  to the identity. Their inverses are admissible harmonic charts on the
  approximating manifolds. The remaining properties follow as before.
\end{proof}


\subsection{Compact Classes of Riemannian Manifolds}

\begin{lemma}[Fundamental Theorem, part A — Gromov--Hausdorff precompactness]
  \label{lem:pet-ch11-fundamental-thm-precompact-gh}
  \lean{PetersenLib.fundamentalTheorem_ghPrecompact}
  \source{petersen:page-0434,petersen:page-0435}
  \uses{def:pet-ch11-cmalpha-norm-pointed, cor:pet-ch11-ricci-diam-precompactness}
  Fix $Q>0$, $n\ge2$, $m\ge0$, $\alpha\in(0,1]$, and $r>0$, and let
  $\mathcal{M}=\mathcal{M}^{m,\alpha}(n,Q,r)$ be the class of complete,
  pointed Riemannian $n$-manifolds $(M,g,p)$ with
  $\|(M,g)\|_{C^{m,\alpha},r}\le Q$. Then $\mathcal{M}$ is precompact in the
  pointed Gromov--Hausdorff topology.
\end{lemma}
\begin{proof}
  \uses{lem:pet-ch11-fundamental-thm-precompact-gh}
  \sketch
  Put \(\delta=e^{-Q}r\). The chart bounds cover each \(\delta\)-ball by a
  fixed number \(N(n,Q)\) of balls of radius \(\delta/4\); induction covers
  \(B(p,R)\) by \(N^\ell\) such balls with \(\ell\delta/2\ge R\). Each
  disjoint \(\varepsilon\)-ball has volume at least
  \(e^{-nQ}\vol B(0,\varepsilon)\), while the preceding cover bounds
  \(\vol B(p,R)\). Comparing total volume with the smallest disjoint ball
  gives a uniform capacity bound. Applying
  \cref{prop:pet-ch11-gromov-precompactness-criterion} on bounded balls and
  then \cref{cor:pet-ch11-pointed-precompactness-criterion} proves pointed GH
  precompactness.
\end{proof}


\begin{lemma}[Fundamental Theorem, part B — the limit is a norm-bounded manifold]
  \label{lem:pet-ch11-fundamental-thm-limit-is-manifold}
  \lean{PetersenLib.fundamentalTheorem_limitIsManifold}
  \source{petersen:page-0435}
  \uses{lem:pet-ch11-fundamental-thm-precompact-gh, prop:pet-ch11-norm-properties}
  For the class in \cref{lem:pet-ch11-fundamental-thm-precompact-gh}, if
  $(M_i,g_i,p_i)\in\mathcal{M}$ converge to $(X,|\cdot|,p)$ in the pointed
  Gromov--Hausdorff topology, then $X$ carries the structure of a $C^{m,\alpha}$
  Riemannian manifold with $\|(X,g)\|_{C^{m,\alpha},r}\le Q$, whose induced
  distance agrees with $|\cdot|$.
\end{lemma}
\begin{proof}
  \uses{lem:pet-ch11-fundamental-thm-limit-is-manifold}
  \sketch
  Approximate each \(q\in X\) by \(q_i\) and choose uniformly controlled
  charts \(\varphi_i\) centered at \(q_i\). Arzelà--Ascoli gives a limiting
  chart \(\varphi\), and the distance bounds in
  \cref{prop:pet-ch11-norm-properties} make it a homeomorphism onto its image.
  The pulled-back metrics converge to \(g_*\) satisfying the same coefficient
  bounds and inducing the limit distance. On overlaps, the limiting
  distance-preserving transitions have the required regularity, producing a
  \(C^{m,\alpha}\) Riemannian structure on \(X\).
\end{proof}


\begin{theorem}[Fundamental Theorem of Convergence Theory]
  \label{thm:pet-ch11-fundamental-theorem-convergence}
  \lean{PetersenLib.fundamentalTheoremOfConvergenceTheory}
  \source{petersen:page-0434,petersen:page-0435,petersen:page-0436,petersen:page-0437}
  \uses{lem:pet-ch11-fundamental-thm-precompact-gh, lem:pet-ch11-fundamental-thm-limit-is-manifold}
  \dcref{ch11:11.3.6}
  Fix $Q>0$, $n\ge2$, $m\ge0$, $\alpha\in(0,1]$, and $r>0$. The class
  $\mathcal{M}^{m,\alpha}(n,Q,r)$ of complete pointed Riemannian $n$-manifolds
  $(M,g,p)$ with $\|(M,g)\|_{C^{m,\alpha},r}\le Q$ is compact in the pointed
  $C^{m,\beta}$ topology for all $\beta<\alpha$.
\end{theorem}
\begin{proof}
  \uses{thm:pet-ch11-fundamental-theorem-convergence}
  \sketch
  By \cref{lem:pet-ch11-fundamental-thm-precompact-gh}, pass to a pointed GH limit; \cref{lem:pet-ch11-fundamental-thm-limit-is-manifold} supplies its \(C^{m,\alpha}\) metric. Choose countable limiting atlases on the limit and the approximants. On each finite union, glue the local comparison maps by a partition of unity in a common chart; convergence of transition maps makes the result \(C^{m+1,\beta}\)-close to every local map. A diagonal exhaustion yields embeddings whose pulled-back metrics converge in \(C^{m,\beta}\), for every \(\beta<\alpha\).
\end{proof}


\begin{corollary}[diameter/volume bounds give finiteness]
  \label{cor:pet-ch11-fundamental-thm-diam-vol-bounded}
  \lean{PetersenLib.fundamentalTheorem_diamOrVolBounded_finiteDiffeoTypes}
  \source{petersen:page-0437}
  \uses{thm:pet-ch11-fundamental-theorem-convergence}
  \dcref{ch11:11.3.7}
  Imposing either $\diam\le D$ or $\vol\le V$ on
  $\mathcal{M}^{m,\alpha}(n,Q,r)$ gives subclasses whose elements satisfy
  $\diam\le D$, respectively $\vol\le V$, and are compact in the
  $C^{m,\beta}$ topology (for $\beta<\alpha$); in particular each such
  subclass contains only finitely many diffeomorphism types.
\end{corollary}
\begin{proof}
  \uses{cor:pet-ch11-fundamental-thm-diam-vol-bounded}
  \sketch
  In \cref{thm:pet-ch11-fundamental-theorem-convergence}, a diameter bound restricts the exhaustion to uniformly many controlled charts, so pointed convergence is unpointed. If \(\vol M\le V\), the local chart-volume lower bound gives a uniform bound
  \[
    V e^{2nQ}\bigl(\vol B(0,\varepsilon)\bigr)^{-1}
  \]
  on disjoint \(\varepsilon\)-balls, hence a diameter bound. Compactness then permits only finitely many diffeomorphism types.
\end{proof}


\begin{corollary}[compactness for the harmonic norm]
  \label{cor:pet-ch11-harmonic-norm-class-compact}
  \lean{PetersenLib.harmonicNormClass_closed_compact}
  \source{petersen:page-0437}
  \uses{prop:pet-ch11-harmonic-norm-properties, thm:pet-ch11-fundamental-theorem-convergence}
  \dcref{ch11:11.3.8}
  Fix $Q>0$, $n\ge2$, $m\ge0$, $\alpha\in(0,1]$, and $r>0$. The class of complete
  pointed Riemannian $n$-manifolds $(M,g,p)$ with
  $\|(M,g)\|^{\mathrm{har}}_{C^{m,\alpha},r}\le Q$ is closed in the pointed
  $C^{m,\alpha}$ topology and compact in the pointed $C^{m,\beta}$ topology
  for all $\beta<\alpha$.
\end{corollary}
\begin{proof}
  \uses{cor:pet-ch11-harmonic-norm-class-compact}
  \sketch
  The convergence properties in \cref{prop:pet-ch11-harmonic-norm-properties} reduce compactness to \cref{thm:pet-ch11-fundamental-theorem-convergence}. Limits of harmonic charts remain harmonic, so the norm bound is closed.
\end{proof}


\subsection{Alternative Norms}

\begin{remark}[alternative norm conditions]
  \label{rem:pet-ch11-alternative-norm-conditions}
  \lean{PetersenLib.alternativeNormConditions}
  \source{petersen:page-0437}
  \uses{def:pet-ch11-cmalpha-norm-pointed}
  Replacing (n1), (n2) by
  \[
    \text{(n1$'$)}\quad|D\varphi|,|D\varphi^{-1}|\le f_1(n,Q),\qquad
    \text{(n2$'$)}\quad r^{|I|+\alpha}\|\partial^Ig\|_\alpha\le f_2(n,Q),\ 0\le|I|\le m,
  \]
  for continuous $f_1,f_2$ with $f_1(n,0)=1$, $f_2(n,0)=0$: this preserves
  continuity of the norm in $r$, the Fundamental Theorem, and the
  characterization of flat manifolds and Euclidean space.
\end{remark}

\begin{remark}[the $C^0$ norm and topological finiteness]
  \label{rem:pet-ch11-c0-norm-precompactness-only}
  \lean{PetersenLib.c0NormClass_precompact_gh_only}
  \source{petersen:page-0438}
  \uses{def:pet-ch11-cmalpha-norm-global, cor:pet-ch11-fundamental-thm-diam-vol-bounded}
  At $m=\alpha=0$, condition (n2) is unavailable, but a $C^0$-norm notion
  survives. The class $\mathcal{M}^0(n,Q,r)$ is only precompact in the
  pointed Gromov--Hausdorff topology (the flat-manifold characterization
  still holds), and diameter- or volume-bounded subclasses are only
  Gromov--Hausdorff precompact, with diffeomorphism-type finiteness
  apparently failing: the averaging construction of
  \cref{thm:pet-ch11-fundamental-theorem-convergence} breaks down since the
  coordinates converge only in $C^{\alpha}$, $\alpha<1$. Using local
  contractibility of homeomorphism groups, two $C^0$-close topological
  embeddings can nonetheless be glued without much $C^0$ distortion, giving
  topological embeddings with converging pullback metrics; hence
  diameter- or volume-bounded classes contain only finitely many
  homeomorphism types.
\end{remark}

\begin{definition}[Reifenberg norm]
  \label{def:pet-ch11-reifenberg-norm}
  \lean{PetersenLib.reifenbergNorm}
  \source{petersen:page-0438}
  \uses{def:pet-ch11-gromov-hausdorff-distance}
  For a metric space $(X,|\cdot|)$, define the \emph{$n$-dimensional Reifenberg norm
  at scale $r$} by
  \[
    \|(X,|\cdot|)\|^n_r=\frac1r\sup_{p\in X}\dGH\big(B(p,r),B(0,r)\big),
  \]
  where $B(0,R)\subset\mathbb{R}^n$; the factor $r^{-1}$ prevents small distances
  merely because $r$ is small. If $(X_i,|\cdot|_i)\to(X,|\cdot|)$ in the
  Gromov--Hausdorff topology then $\|(X_i,|\cdot|_i)\|^n_r\to\|(X,|\cdot|)\|^n_r$
  for fixed $n,r$.
\end{definition}

\begin{remark}[Reifenberg norm: vanishing and the Cheeger--Colding converse]
  \label{rem:pet-ch11-reifenberg-norm-properties}
  \lean{PetersenLib.reifenbergNorm_limit_zero,PetersenLib.cheegerColdingReifenbergConverse}
  \source{petersen:page-0438,petersen:page-0439}
  \uses{def:pet-ch11-reifenberg-norm}
  Every $n$-dimensional Riemannian manifold satisfies
  $\lim_{r\to0}\|(M,g)\|^n_r=0$. Cheeger and Colding proved a converse: there
  is $\varepsilon(n)>0$ such that if $\|(X,|\cdot|)\|^n_r\le\varepsilon(n)$
  for all small $r$, then $X$ is, in a weak sense, an $n$-dimensional
  Riemannian manifold; in particular, for small $r$, the $\alpha$-Hölder
  distance between $B(p,r)$ and $B(0,r)$ is then small.
\end{remark}

\begin{definition}[$\alpha$-Hölder distance]
  \label{def:pet-ch11-holder-distance}
  \lean{PetersenLib.holderDistance}
  \source{petersen:page-0438,petersen:page-0439}
  \uses{def:pet-ch11-holder-pseudonorm}
  For metric spaces $X,Y$, define their \emph{$\alpha$-Hölder distance}
  $d_\alpha(X,Y)$ as
  the infimum, over homeomorphisms $F:X\to Y$, of
  \[
    \log\max\Big\{\sup_{x_1\ne x_2}\frac{|F(x_1)F(x_2)|}{|x_1x_2|^\alpha},\
    \sup_{y_1\ne y_2}\frac{|F^{-1}(y_1)F^{-1}(y_2)|}{|y_1y_2|^\alpha}\Big\}.
  \]
\end{definition}

\begin{remark}[Cheeger--Colding Hölder convergence and Colding's volume criterion]
  \label{rem:pet-ch11-cheeger-colding-holder-convergence}
  \lean{PetersenLib.cheegerColdingHolderConvergence,PetersenLib.coldingVolumeCriterion}
  \source{petersen:page-0439}
  \uses{def:pet-ch11-holder-distance, rem:pet-ch11-reifenberg-norm-properties}
  Under Gromov--Hausdorff convergence $(M_i,g_i)\to(X,|\cdot|)$, Cheeger and
  Colding show that if $\|(M_i,g_i)\|^n_r\le\varepsilon(n)$ for all
  $i$ and small $r$, then $(M_i,g_i)\to(X,|\cdot|)$ in the Hölder distance,
  and in particular all $M_i$ are diffeomorphic to $X$ for large $i$.
  Colding earlier showed that for $(M,g)$ with $\Ric\ge(n-1)k$,
  $\|(M,g)\|^n_r$ is small iff $\vol B(p,r)\ge(1-\delta)\vol B(0,r)$ for some
  small $\delta$; relative volume comparison shows the volume condition for
  one $r$ persists for all smaller $r$, so smallness of the Reifenberg norm
  for all small $r$ follows from the volume condition at a single $r$.
\end{remark}

\section{Geometric Applications}

\subsection{Ricci Curvature}

\begin{lemma}[Ricci controls the harmonic $C^{1,\alpha}$ norm]
  \label{lem:pet-ch11-anderson-ricci-harmonic-c1alpha-bound}
  \lean{PetersenLib.andersonRicciHarmonicC1AlphaBound}
  \source{petersen:page-0439,petersen:page-0440}
  \uses{def:pet-ch11-harmonic-norm, lem:pet-ch11-laplacian-ricci-harmonic-coords, thm:pet-ch11-schauder-estimates}
  \dcref{ch11:11.4.1}
  If $|\Ric|\le\Lambda$ and $r_1<r_2$,
  $K\ge\|(M,g,p)\|^{\mathrm{har}}_{C^{1,r_2}}$, $\alpha\in(0,1)$ there is
  $C(n,\alpha,K,r_1,r_2,\Lambda)$ with
  $\|(M,g,p)\|^{\mathrm{har}}_{C^{1,\alpha},r_1}\le C$. If moreover $g$ is
  Einstein, $\Ric=kg$, then for each $m$ there is
  $C(n,\alpha,K,r_1,r_2,k,m)$ with
  $\|(M,g,p)\|^{\mathrm{har}}_{C^{m+1,\alpha},r_1}\le C$.
\end{lemma}
\begin{proof}
  \uses{lem:pet-ch11-anderson-ricci-harmonic-c1alpha-bound}
  \sketch
  In a controlled harmonic chart, apply \cref{thm:pet-ch11-schauder-estimates} to the metric equation from \cref{lem:pet-ch11-laplacian-ricci-harmonic-coords},
  \[
    \Delta g_{ij}=-2\Ric_{ij}-2Q(g,\partial g).
  \]
  The \(C^1\) chart bound controls ellipticity and the quadratic term, while \(|\Ric|\le\Lambda\) controls the remainder, giving the \(C^{1,\alpha}\) estimate on the smaller ball. If \(\Ric=kg\), the right side gains the regularity of \(g\); iterated Schauder estimates yield every \(C^{m+1,\alpha}\) bound.
\end{proof}


\begin{corollary}[precompactness under bounded Ricci and harmonic norm]
  \label{cor:pet-ch11-ricci-bounded-precompact-c1alpha}
  \lean{PetersenLib.ricciAndHarmonicNormBounded_precompact}
  \source{petersen:page-0441}
  \uses{lem:pet-ch11-anderson-ricci-harmonic-c1alpha-bound, thm:pet-ch11-fundamental-theorem-convergence}
  \dcref{ch11:11.4.2}
  With $n\ge2$ and $Q,r,\Lambda\in(0,\infty)$, the class of Riemannian
  $n$-manifolds with $\|(M,g)\|^{\mathrm{har}}_{C^{1,r}}\le Q$,
  $|\Ric|\le\Lambda$ is precompact in the pointed $C^{1,\alpha}$ topology for
  every $\alpha\in(0,1)$; the subclass of Einstein manifolds is compact in
  the $C^{m,\alpha}$ topology for every $m\ge1$, $\alpha\in(0,1)$.
\end{corollary}

\begin{theorem}[injectivity radius controls the harmonic norm]
  \label{thm:pet-ch11-anderson-injectivity-radius-controls-norm}
  \lean{PetersenLib.andersonInjectivityRadiusControlsNorm}
  \source{petersen:page-0441}
  \uses{lem:pet-ch11-anderson-ricci-harmonic-c1alpha-bound, thm:pet-ch11-fundamental-theorem-convergence, rem:pet-ch11-einstein-harmonic-regularity-bootstrap}
  \dcref{ch11:11.4.3}
  For $n\ge2$, $\alpha\in(0,1)$, and $\Lambda,R>0$, each $Q>0$ admits
  $r(n,\alpha,\Lambda,R)>0$ such that any compact Riemannian $n$-manifold
  $(M,g)$ with $|\Ric|\le\Lambda$, $\inj\ge R$ satisfies
  $\|(M,g)\|^{\mathrm{har}}_{C^{1,\alpha},r}\le Q$.
\end{theorem}
\begin{proof}
  \uses{thm:pet-ch11-anderson-injectivity-radius-controls-norm}
  \sketch
  Otherwise choose scales \(r_i\downarrow0\) at which the harmonic norm equals a fixed \(Q\), and rescale by \(r_i^{-2}\). Then Ricci tends to zero, injectivity radius tends to infinity, and the unit-scale norm remains in \([Q/2,Q]\) at suitable base points. The estimate in \cref{lem:pet-ch11-anderson-ricci-harmonic-c1alpha-bound} and \cref{thm:pet-ch11-fundamental-theorem-convergence} give a \(C^{1,\gamma}\) limit. The equation in \cref{lem:pet-ch11-laplacian-ricci-harmonic-coords}, followed by the regularity in \cref{rem:pet-ch11-einstein-harmonic-regularity-bootstrap}, makes the limit smooth and Ricci flat. Infinite injectivity radius and splitting force it to be Euclidean, contradicting the nonzero limiting norm.
\end{proof}


\begin{corollary}[diffeomorphism finiteness under Ricci, diameter, injectivity bounds]
  \label{cor:pet-ch11-anderson-diffeo-finiteness-injectivity}
  \lean{PetersenLib.andersonDiffeoFinitenessInjectivity}
  \source{petersen:page-0442,petersen:page-0443}
  \uses{thm:pet-ch11-anderson-injectivity-radius-controls-norm, cor:pet-ch11-fundamental-thm-diam-vol-bounded}
  \dcref{ch11:11.4.4}
  Closed Riemannian $n$-manifolds with $n\ge2$ and
  $\Lambda,D,R>0$ satisfying
  $|\Ric|\le\Lambda$, $\diam\le D$, $\inj\ge R$ is precompact in the
  $C^{1,\alpha}$ topology for every $\alpha\in(0,1)$, and in particular
  contains only finitely many diffeomorphism types.
\end{corollary}

\subsection{Volume Pinching}

\begin{lemma}[near-maximal volume growth forces Euclidean space]
  \label{lem:pet-ch11-anderson-volume-pinching-euclidean}
  \lean{PetersenLib.andersonVolumePinchingEuclidean}
  \source{petersen:page-0443,petersen:page-0444}
  \uses{thm:pet-ch11-fundamental-theorem-convergence, cor:pet-ch11-norm-zero-implies-flat}
  \dcref{ch11:11.4.5}
  For every $n\ge2$ one can choose $\varepsilon(n)>0$ so that a complete Ricci
  flat $(M,g)$ with $\vol B(p,r)\ge(1-\varepsilon)\omega_nr^n$ for some
  $p\in M$ and some $r$ is isometric to Euclidean space.
\end{lemma}
\begin{proof}
  \uses{lem:pet-ch11-anderson-volume-pinching-euclidean}
  \sketch
  Bishop--Gromov monotonicity makes the asymptotic volume ratio independent of the base point and invariant under scaling. If no \(\varepsilon(n)\) worked, rescale counterexamples so their unit-scale norm lies in \([1/2,1]\). The compactness theorem \cref{thm:pet-ch11-fundamental-theorem-convergence} gives a Ricci-flat limit with Euclidean volume ratio \(1\). Equality in volume comparison makes every ball Euclidean, hence the limit norm vanishes by \cref{cor:pet-ch11-norm-zero-implies-flat}, contradicting continuity of the normalized norm.
\end{proof}


\begin{corollary}[volume-pinching precompactness]
  \label{cor:pet-ch11-volume-pinching-precompactness}
  \lean{PetersenLib.volumePinchingPrecompactness}
  \source{petersen:page-0444,petersen:page-0445}
  \uses{lem:pet-ch11-anderson-volume-pinching-euclidean, thm:pet-ch11-fundamental-theorem-convergence}
  \dcref{ch11:11.4.6}
  Given $n\ge2$, $-\infty<\lambda\le\Lambda<\infty$, and $D,R>0$, choose
  $\delta=\delta(n,\lambda R^2)$ so that the class of closed Riemannian
  $n$-manifolds with $(n-1)\Lambda\ge\Ric\ge(n-1)\lambda$, $\diam\le D$,
  $\vol B(p,R)\ge(1-\delta)v(n,\lambda,R)$ is precompact in the $C^{1,\alpha}$
  topology for every $\alpha\in(0,1)$, and in particular contains only
  finitely many diffeomorphism types.
\end{corollary}
\begin{proof}
  \uses{cor:pet-ch11-volume-pinching-precompactness}
  \sketch
  Use the blow-up normalization from \cref{thm:pet-ch11-anderson-injectivity-radius-controls-norm}. Under \(\bar g_i=k_i^2g_i\), relative volume comparison gives, for fixed \(r\),
  \[
    \vol_{\bar g_i}B(p_i,r)\ge(1-\delta)v(n,\lambda k_i^{-2},r)
      \longrightarrow(1-\delta)\omega_nr^n .
  \]
  The compactness theorem \cref{thm:pet-ch11-fundamental-theorem-convergence} produces a Ricci-flat limit, which is Euclidean by \cref{lem:pet-ch11-anderson-volume-pinching-euclidean}. This contradicts the normalized nonzero harmonic norm and proves the required uniform control.
\end{proof}


\subsection{Sectional Curvature}

\begin{theorem}[the Convergence Theorem of Riemannian Geometry]
  \label{thm:pet-ch11-convergence-theorem-riemannian-geometry}
  \lean{PetersenLib.convergenceTheoremOfRiemannianGeometry}
  \source{petersen:page-0445}
  \uses{thm:pet-ch11-anderson-injectivity-radius-controls-norm}
  \dcref{ch11:11.4.7}
  For $R,K>0$, there are $Q,r>0$ such that every $(M,g)$ with $\inj\ge R$
  and $|\sec|\le K$ satisfies
  $\|(M,g)\|^{\mathrm{har}}_{C^{1,\alpha},r}\le Q$; in
  particular this class is compact in the pointed $C^{1,\alpha}$ topology for
  all $\alpha<1$.
\end{theorem}

\begin{corollary}[even-dimensional pinched curvature finiteness]
  \label{cor:pet-ch11-cheeger-even-dim-finiteness}
  \lean{PetersenLib.cheegerEvenDimFiniteness}
  \source{petersen:page-0445}
  \uses{thm:pet-ch11-convergence-theorem-riemannian-geometry}
  \dcref{ch11:11.4.8}
  For $n\ge1$ and $k>0$, the class of Riemannian $2n$-manifolds satisfying
  $k\le\sec\le1$ is compact in the $C^{1,\alpha}$ topology and hence contains
  only finitely many diffeomorphism types (using the diameter bound in
  positive curvature and Klingenberg's injectivity-radius estimate).
\end{corollary}

\begin{lemma}[volume lower bound controls injectivity radius]
  \label{lem:pet-ch11-cheeger-injectivity-radius-lemma}
  \lean{PetersenLib.cheegerInjectivityRadiusLemma}
  \source{petersen:page-0446}
  \uses{thm:pet-ch11-fundamental-theorem-convergence, cor:pet-ch11-norm-zero-implies-flat}
  \dcref{ch11:11.4.9}
  If $n\ge2$, $v,K>0$, and compact $(M,g)$ obeys
  $|\sec|\le K$ and $\vol B(p,1)\ge v$ for all $p\in M$, then $\inj M\ge R$
  for $R=R(n,K,v)$.
\end{lemma}
\begin{proof}
  \uses{lem:pet-ch11-cheeger-injectivity-radius-lemma}
  \sketch
  Rescale a sequence with \(\inj M_i\to0\) by \((\inj M_i)^{-2}\). The rescaled spaces have injectivity radius \(1\), curvature tending to zero, and by \cref{thm:pet-ch11-fundamental-theorem-convergence} converge to a flat pointed limit. Klingenberg gives length-\(2\) geodesic loops, so the limit has finite injectivity radius. The volume hypothesis and relative comparison yield Euclidean-order volume growth at every scale. A nontrivial complete flat quotient has at most \((n-1)\)-dimensional growth after passage to a cyclic cover, so the limit is \(\mathbb R^n\), whose injectivity radius is infinite: a contradiction.
\end{proof}


\begin{corollary}[bounded curvature, diameter, volume finiteness]
  \label{cor:pet-ch11-cheeger-diam-vol-finiteness}
  \lean{PetersenLib.cheegerDiamVolFiniteness}
  \source{petersen:page-0447}
  \uses{lem:pet-ch11-cheeger-injectivity-radius-lemma, thm:pet-ch11-convergence-theorem-riemannian-geometry}
  \dcref{ch11:11.4.10}
  With $n\ge2$ and $K,D,v>0$, closed Riemannian $n$-manifolds satisfying
  $|\sec|\le K$, $\diam\le D$, $\vol\ge v$ is precompact in the $C^{1,\alpha}$
  topology for every $\alpha\in(0,1)$, and in particular contains only
  finitely many diffeomorphism types.
\end{corollary}

\subsection{Lower Curvature Bounds}

\begin{theorem}[lower sectional curvature and injectivity radius control the norm]
  \label{thm:pet-ch11-lower-sec-bound-injectivity-norm-control}
  \lean{PetersenLib.lowerSecBoundInjectivityNormControl}
  \source{petersen:page-0447}
  \uses{def:pet-ch11-cmalpha-norm-pointed}
  \dcref{ch11:11.4.11}
  For $R,k>0$, there are $Q,r$ depending on $R,k$ such that every $(M,g)$
  with $\sec\ge-k^2$, $\inj\ge R$ satisfies $\|(M,g)\|_{C^{1,r}}\le Q$.
\end{theorem}
\begin{proof}
  \uses{thm:pet-ch11-lower-sec-bound-injectivity-norm-control}
  \sketch
  Hessian comparison and the lower support \(R-|x\,c(R)|\) give uniform two-sided Hessian bounds for distance functions away from their centers. Fixing an orthonormal basis at \(p\), use
  \[
    x\longmapsto\bigl(|x\,c_1(-R/2)|,\ldots,|x\,c_n(-R/2)|\bigr)
  \]
  as coordinates near \(p\). Their inverse metric coefficients are inner products of the corresponding unit gradients; the Hessian bounds control their first derivatives. On a uniform smaller ball the coefficient matrix stays invertible and satisfies (n1)--(n2), giving the asserted \(Q,r\).
\end{proof}


\begin{example}[optimality: only $C^{1,1}$ limits under lower curvature]
  \label{ex:pet-ch11-optimality-lower-curvature-example}
  \lean{PetersenLib.optimalityLowerCurvatureExample}
  \source{petersen:page-0448}
  \uses{thm:pet-ch11-lower-sec-bound-injectivity-norm-control}
  \dcref{ch11:11.4.12}
  As $\varepsilon\to0$, consider rotationally symmetric metrics
  $dr^2+\phi_\varepsilon^2(r)d\theta^2$ with $\phi_\varepsilon$ concave
  $\phi_\varepsilon(r)=r$ for $0\le r\le1-\varepsilon$ and
  $\phi_\varepsilon(r)=\tfrac34r$ for $r\ge1+\varepsilon$, have $\sec\ge0$,
  $\inj\ge1$, and converge to a $C^{1,1}$ manifold with $C^{0,1}$ metric
  $(M,g)$, for which $\|(M,g)\|_{C^{1,r}}<\infty$ for all $r$: limit spaces
  of sequences with $\inj\ge R$, $\sec\ge-k^2$ cannot in general be smoother
  than this.
\end{example}

\begin{example}[a $C^{1,1}$ limit with two-sided curvature bound]
  \label{ex:pet-ch11-c11-limit-example}
  \lean{PetersenLib.c11LimitExample}
  \source{petersen:page-0448}
  \uses{ex:pet-ch11-optimality-lower-curvature-example}
  \dcref{ch11:11.4.13}
  Take $\psi_\varepsilon(r)=\sin r$ for $0\le r\le\pi/2-\varepsilon$ and
  $\psi_\varepsilon(r)=1$ for $r\ge\pi/2$; the metrics
  $dr^2+\psi_\varepsilon^2(r)d\theta^2$ satisfy $|\sec|\le4$, $\inj\ge1/4$; as
  $\varepsilon\to0$ they converge to a $C^{1,1}$ limit metric, though only
  $C^{1,\alpha}$ convergence for $\alpha<1$ has been established for such
  limits. Unlike the two-sided case, no injectivity-radius lower bound need
  hold under a one-sided curvature bound alone.
\end{example}

\begin{remark}[concave cone smoothing with collapsing injectivity radius]
  \label{rem:pet-ch11-exercise-11-4-14}
  \lean{PetersenLib.exercise_11_4_14}
  \source{petersen:page-0449}
  \uses{ex:pet-ch11-optimality-lower-curvature-example}
  \dcref{ch11:11.4.14}
  For $a\in(0,1)$ and $\varepsilon>0$, choose a smooth concave
  $\rho_\varepsilon(r)$ with $\rho_\varepsilon(r)=r$ for $0\le r\le\varepsilon$
  and $\rho_\varepsilon(r)=ar$ for $r\ge2\varepsilon$. The surfaces
  $dr^2+\rho_\varepsilon^2(r)d\theta^2$ have $\sec\ge0$, $\inj\le5\varepsilon$,
  and the volume of any $R$-ball is always $\ge a\pi R^2$.
\end{remark}

\begin{theorem}[Ricci lower bound with injectivity control]
  \label{thm:pet-ch11-anderson-cheeger-ricci-lower-bound-norm}
  \lean{PetersenLib.andersonCheegerRicciLowerBoundNorm}
  \source{petersen:page-0449}
  \uses{def:pet-ch11-harmonic-norm}
  \dcref{ch11:11.4.15}
  For $R,k>0$ and $\alpha\in(0,1)$, there are suitable $Q,r$ depending on
  $n,R,k$
  such that any $(M^n,g)$ with $\Ric\ge-(n-1)k^2$, $\inj\ge R$ satisfies
  $\|(M,g)\|^{\mathrm{har}}_{C^{\alpha},r}\le Q$.
\end{theorem}

\subsection{Curvature Pinching}

\begin{theorem}[Ricci pinching implies $C^{1,\alpha}$-closeness to Einstein]
  \label{thm:pet-ch11-ricci-pinching-einstein-closeness}
  \lean{PetersenLib.ricciPinchingEinsteinCloseness}
  \source{petersen:page-0450}
  \uses{cor:pet-ch11-anderson-diffeo-finiteness-injectivity, rem:pet-ch11-einstein-harmonic-regularity-bootstrap}
  \dcref{ch11:11.4.16}
  For $n\ge2$, $R,D>0$, and $\lambda\in\mathbb{R}$, there is
  $\varepsilon(n,\lambda,D,R)>0$ such that any closed Riemannian $n$-manifold
  $(M,g)$ with $\diam\le D$, $\inj\ge R$, $|\Ric-\lambda g|\le\varepsilon$ is
  $C^{1,\alpha}$ close to an Einstein metric with Einstein constant $\lambda$.
\end{theorem}
\begin{proof}
  \uses{thm:pet-ch11-ricci-pinching-einstein-closeness}
  \sketch
  The class is \(C^{1,\alpha}\)-precompact by \cref{cor:pet-ch11-anderson-diffeo-finiteness-injectivity}. A counterexample sequence therefore converges to \(g\), and the harmonic-coordinate equation in \cref{lem:pet-ch11-laplacian-ricci-harmonic-coords} passes to the weak limit:
  \[
    \tfrac12\Delta g+Q(g,\partial g)=-\lambda g.
  \]
  The bootstrap in \cref{rem:pet-ch11-einstein-harmonic-regularity-bootstrap} makes \(g\) smooth and Einstein, contradicting the assumed failure of closeness.
\end{proof}


\begin{remark}[injectivity radius may be replaced by almost-maximal volume]
  \label{rem:pet-ch11-ricci-pinching-volume-alternative}
  \lean{PetersenLib.ricciPinchingVolumeAlternative}
  \source{petersen:page-0450}
  \uses{thm:pet-ch11-ricci-pinching-einstein-closeness, cor:pet-ch11-volume-pinching-precompactness}
  Replacing the injectivity-radius argument by
  \cref{cor:pet-ch11-volume-pinching-precompactness} in the
  proof of
  \cref{thm:pet-ch11-ricci-pinching-einstein-closeness}, the injectivity
  radius hypothesis there may be replaced by an almost-maximal-volume
  hypothesis.
\end{remark}

\begin{theorem}[sectional-curvature pinching implies closeness to constant curvature]
  \label{thm:pet-ch11-sec-pinching-constant-curvature-closeness}
  \lean{PetersenLib.secPinchingConstantCurvatureCloseness}
  \source{petersen:page-0450}
  \uses{thm:pet-ch11-ricci-pinching-einstein-closeness, lem:pet-ch11-cheeger-injectivity-radius-lemma}
  \dcref{ch11:11.4.17}
  For $n\ge2$, $v,D>0$, and $\lambda\in\mathbb{R}$, there is
  $\varepsilon(n,\lambda,D,v)>0$ such that any closed Riemannian $n$-manifold
  $(M,g)$ with $\diam\le D$, $\vol\ge v$, $|\sec-\lambda|\le\varepsilon$ is
  $C^{1,\alpha}$ close to a metric of constant curvature $\lambda$.
\end{theorem}
\begin{proof}
  \uses{thm:pet-ch11-sec-pinching-constant-curvature-closeness}
  \sketch
  By \cref{lem:pet-ch11-cheeger-injectivity-radius-lemma}, curvature, diameter, and volume bounds give a uniform injectivity radius; \cref{thm:pet-ch11-ricci-pinching-einstein-closeness} then gives an Einstein limit with \(\Ric=(n-1)\lambda g\). In polar coordinates, Hessian comparison for \(\lambda-\varepsilon_i\le\sec_{g_i}\le\lambda+\varepsilon_i\) passes to
  \[
    \Hess r=\frac{\mathrm{sn}_{\lambda}^{\prime}(r)}{\mathrm{sn}_{\lambda}(r)}g_r.
  \]
  The equality case forces constant sectional curvature \(\lambda\).
\end{proof}


\begin{corollary}[even-dimensional positive pinching without hypotheses]
  \label{cor:pet-ch11-even-dim-positive-pinching}
  \lean{PetersenLib.evenDimPositivePinching}
  \source{petersen:page-0451}
  \uses{thm:pet-ch11-sec-pinching-constant-curvature-closeness}
  \dcref{ch11:11.4.18}
  In even dimension $2n\ge2$ and for $\lambda>0$, choose
  $\varepsilon=\varepsilon(n,\lambda)>0$
  such that any closed Riemannian $2n$-manifold $(M,g)$ with
  $|\sec-\lambda|\le\varepsilon$ is $C^{1,\alpha}$ close to a metric of
  constant curvature $\lambda$ (in even dimensions, positive curvature alone
  bounds the diameter and, together with an upper curvature bound, the
  injectivity radius from below).
\end{corollary}

\begin{corollary}[simply connected positive pinching]
  \label{cor:pet-ch11-odd-dim-simply-connected-pinching}
  \lean{PetersenLib.oddDimSimplyConnectedPinching}
  \source{petersen:page-0451}
  \uses{cor:pet-ch11-even-dim-positive-pinching}
  \dcref{ch11:11.4.19}
  For $n\ge2$ and $\lambda>0$, choose $\varepsilon=\varepsilon(n,\lambda)>0$
  such that any closed simply connected Riemannian $n$-manifold $(M,g)$ with
  $|\sec-\lambda|\le\varepsilon$ is $C^{1,\alpha}$ close to a metric of
  constant curvature $\lambda$.
\end{corollary}

\section{Exercises}

\begin{remark}[low-dimensional Hausdorff approximations of a cube]
  \label{rem:pet-ch11-exercise-11-6-1}
  \lean{PetersenLib.exercise_11_6_1}
  \source{petersen:page-0453}
  \uses{def:pet-ch11-hausdorff-distance}
  \dcref{ch11:11.6.1}
  Construct a sequence of $1$-dimensional metric spaces converging in Hausdorff
  distance to the unit cube $[0,1]^3$ with the metric from the maximum norm on
  $\mathbb{R}^3$. Then find surfaces (jungle gyms) converging to the same
  space.
\end{remark}

\begin{remark}[\(\varepsilon\)-isometries control Gromov--Hausdorff distance]
  \label{rem:pet-ch11-exercise-11-6-2}
  \lean{PetersenLib.exercise_11_6_2}
  \source{petersen:page-0453}
  \uses{def:pet-ch11-gromov-hausdorff-distance}
  \dcref{ch11:11.6.2}
  Suppose a (not necessarily continuous) map $F:X\to Y$ satisfies
  $\big||x_1x_2|-|F(x_1)F(x_2)|\big|\le\varepsilon$ for all $x_1,x_2\in X$ and
  $F(X)\subset Y$ is $\varepsilon$-dense, show $\dGH(X,Y)<2\varepsilon$.
\end{remark}

\begin{remark}[injectivity radius gives local volume and precompactness]
  \label{rem:pet-ch11-exercise-11-6-3}
  \lean{PetersenLib.exercise_11_6_3}
  \source{petersen:page-0453}
  \uses{prop:pet-ch11-gromov-precompactness-criterion}
  \dcref{ch11:11.6.3}
  Croke's estimate gives a universal $c(n)$ for which every $n$-manifold
  with $\inj\ge R$ satisfies $\vol B(p,r)\ge c(n)r^n$ for $r\le R/e$; use this
  to show the class of $n$-manifolds with $\inj\ge R$, $\vol\le V$ is
  precompact in the Gromov--Hausdorff topology.
\end{remark}

\begin{remark}[bounded-overlap covers under a Ricci lower bound]
  \label{rem:pet-ch11-exercise-11-6-4}
  \lean{PetersenLib.exercise_11_6_4}
  \source{petersen:page-0453}
  \uses{cor:pet-ch11-ricci-diam-precompactness}
  \dcref{ch11:11.6.4}
  For complete $(M,g)$ with $\Ric\ge(n-1)k$, establish a constant
  $C(n,k)$ such that for each $\varepsilon\in(0,1)$ there is a cover by balls
  $B(x_i,\varepsilon)$ such that no more than $C(n,k)$ of the balls
  $B(x_i,5\varepsilon)$ have nonempty pairwise intersection.
\end{remark}

\begin{remark}[Bochner formulas for symmetric covariant derivatives]
  \label{rem:pet-ch11-exercise-11-6-5}
  \lean{PetersenLib.exercise_11_6_5}
  \source{petersen:page-0453}
  \uses{lem:pet-ch11-laplacian-ricci-harmonic-coords}
  \dcref{ch11:11.6.5}
  Derive Bochner formulas for $\Hess(\tfrac12g(X,Y))$ and
  $\Delta g(X,Y)$, for vector fields $X,Y$ with symmetric $\nabla X,\nabla Y$,
  and use these to prove the formulas of
  \cref{lem:pet-ch11-laplacian-ricci-harmonic-coords} relating Ricci
  curvature to the metric in harmonic coordinates.
\end{remark}

\begin{remark}[failure of a \(C^{1,\gamma}\) vector-field estimate]
  \label{rem:pet-ch11-exercise-11-6-6}
  \lean{PetersenLib.exercise_11_6_6}
  \source{petersen:page-0453}
  \uses{thm:pet-ch11-schauder-estimates}
  \dcref{ch11:11.6.6}
  Demonstrate that, unlike the elliptic estimates of
  \cref{thm:pet-ch11-schauder-estimates}, it is not possible to find
  $C^{1,\gamma}$ bounds for a vector field $X$ in terms of $C^0$ bounds on
  $X$ and $\operatorname{div}X$.
\end{remark}

\begin{remark}[convergence of incomplete manifolds]
  \label{rem:pet-ch11-exercise-11-6-7}
  \lean{PetersenLib.exercise_11_6_7}
  \source{petersen:page-0453}
  \uses{def:pet-ch11-tensor-convergence-compact-subset, cor:pet-ch11-anderson-diffeo-finiteness-injectivity}
  \dcref{ch11:11.6.7}
  Extend the definition of $C^{m,\alpha}$ convergence to incomplete manifolds,
  taking the boundary $\partial$ to be the points in the completion but not
  the manifold.
  Show that the class of incomplete spaces with $|\Ric|\le\Lambda$ and
  $\inj(p)\ge\min\{R,R\cdot d(p,\partial)\}$, $R<1$, is precompact in the
  $C^{1,\alpha}$ topology.
\end{remark}

\begin{remark}[weighted norms and convergence]
  \label{rem:pet-ch11-exercise-11-6-8}
  \lean{PetersenLib.exercise_11_6_8}
  \source{petersen:page-0453}
  \uses{thm:pet-ch11-fundamental-theorem-convergence}
  \dcref{ch11:11.6.8}
  Introduce a \emph{weighted norm} by fixing $\rho(R)>0$ and requiring that, in a
  pointed manifold $(M,g,p)$, the points of $S(p,R)$ have norm $\le\rho(R)$.
  Prove the corresponding fundamental theorem.
\end{remark}

\begin{remark}[uniform small-scale norms on compact classes]
  \label{rem:pet-ch11-exercise-11-6-9}
  \lean{PetersenLib.exercise_11_6_9}
  \source{petersen:page-0453}
  \uses{def:pet-ch11-cmalpha-norm-global}
  \dcref{ch11:11.6.9}
  For a class $\mathcal{M}$ of compact Riemannian $n$-manifolds compact in
  the $C^{m,\alpha}$ topology, show there is $f(r)\to0$ as $r\to0$, depending
  on $\mathcal{M}$, with $\|(M,g)\|_{C^{m,\alpha},r}\le f(r)$ for all
  $M\in\mathcal{M}$.
\end{remark}

\begin{remark}[local models under rescaling]
  \label{rem:pet-ch11-exercise-11-6-10}
  \lean{PetersenLib.exercise_11_6_10}
  \source{petersen:page-0454}
  \uses{def:pet-ch11-cmalpha-norm-global, thm:pet-ch11-fundamental-theorem-convergence}
  \dcref{ch11:11.6.10}
  Define the \emph{local models} of a Riemannian class to be limits
  obtained by rescaling elements of the class by constants $\to\infty$; e.g.\
  $|\sec|\le K$ rescales to $|\sec|\le\lambda^{-2}K\to0$, so the local models
  are the flat manifolds, and adding $\inj\ge R$ forces $\inj\to\infty$ under
  rescaling, so the only local model is Euclidean space. Show that a class of
  spaces is compact in the $C^{m,\alpha}$ topology if every sequence in the
  class, rescaled by constants $\to\infty$, subconverges in $C^{m,\alpha}$ to
  Euclidean space.
\end{remark}

\begin{remark}[nonpositive smoothing of a conical metric]
  \label{rem:pet-ch11-exercise-11-6-11}
  \lean{PetersenLib.exercise_11_6_11}
  \source{petersen:page-0454}
  \uses{def:pet-ch11-gromov-hausdorff-distance}
  \dcref{ch11:11.6.11}
  Starting with the singular metric $dt^2+(at)^2d\theta^2$, $a>1$, on
  $\mathbb{R}^2$, construct a sequence of rotationally symmetric metrics on
  $\mathbb{R}^2$ with
  $\sec\le0$, $\inj=\infty$ converging to it in the Gromov--Hausdorff
  topology.
\end{remark}

\begin{remark}[compactness from Ricci derivative bounds]
  \label{rem:pet-ch11-exercise-11-6-12}
  \lean{PetersenLib.exercise_11_6_12}
  \source{petersen:page-0454}
  \uses{thm:pet-ch11-fundamental-theorem-convergence}
  \dcref{ch11:11.6.12}
  Prove that the class of spaces satisfying $\inj\ge R$ and
  $|\nabla^k\Ric|\le\Lambda$ for $k=0,\dots,m$ is compact in the
  $C^{m+1,\alpha}$ topology.
\end{remark}

\begin{remark}[compactness from conjugate-radius control]
  \label{rem:pet-ch11-exercise-11-6-13}
  \lean{PetersenLib.exercise_11_6_13}
  \source{petersen:page-0454}
  \uses{lem:pet-ch11-cheeger-injectivity-radius-lemma}
  \dcref{ch11:11.6.13}
  For complete or compact $n$-manifolds with
  conjugate radius $\ge R$, $|\Ric|\le\Lambda$, $\vol B(p,1)\ge v$, use the
  techniques of \cref{lem:pet-ch11-cheeger-injectivity-radius-lemma} to show
  a lower injectivity-radius bound, hence compactness in $C^{1,\alpha}$.
\end{remark}

\begin{remark}[Eguchi--Hanson obstruction to Ricci compactness]
  \label{rem:pet-ch11-exercise-11-6-14}
  \lean{PetersenLib.exercise_11_6_14}
  \source{petersen:page-0454}
  \uses{rem:pet-ch11-exercise-11-6-13}
  \dcref{ch11:11.6.14}
  The Eguchi--Hanson metrics show that compactness fails in general for the
  class $|\Ric|\le\Lambda$, $\vol B(p,1)\ge v$ alone; either a
  large lower volume bound (as before) or a lower conjugate-radius bound must
  be assumed.
\end{remark}

\begin{remark}[weak harmonic norms and compactness]
  \label{rem:pet-ch11-exercise-11-6-15}
  \lean{PetersenLib.exercise_11_6_15}
  \source{petersen:page-0454,petersen:page-0455}
  \uses{def:pet-ch11-harmonic-norm}
  \dcref{ch11:11.6.15}
  Define the \emph{weak (harmonic) norm} $\|(M,g)\|^{\mathrm{weak}}_{C^{m,\alpha},r}$
  as in \cref{def:pet-ch11-harmonic-norm}, but require only that the charts
  $\varphi_s:B(0,r)\to U_s$ are immersions, so the
  inverse is only locally defined but still makes sense to be harmonic.
  \begin{enumerate}
  \item Show that if $(M,g)$ has bounded sectional curvature, then for every
    $Q>0$ there is $r>0$ with $\|(M,g)\|^{\mathrm{weak}}_{C^{m,\alpha},r}\le Q$;
    thus the weak norm is a generalized curvature quantity.
  \item Show the class of manifolds with bounded weak norm is precompact in
    the Gromov--Hausdorff topology.
  \item Show $(M,g)$ is flat iff the weak norm vanishes on all scales.
  \end{enumerate}
\end{remark}
